% !TeX program = pdflatex
\documentclass[12pt,a4paper]{article}

\usepackage[T2A]{fontenc}
\usepackage[utf8]{inputenc}
\usepackage[english,russian]{babel}
\usepackage{geometry}
\geometry{left=22mm,right=18mm,top=18mm,bottom=19mm}
\usepackage{setspace}
\setstretch{1.12}
\usepackage{amsmath,amssymb}
\usepackage{booktabs}
\usepackage{array}
\usepackage{longtable}
\usepackage{tabularx}
\usepackage{ragged2e}
\usepackage{hyperref}
\hypersetup{colorlinks=true,linkcolor=black,urlcolor=blue}
\usepackage{enumitem}
\setlist{itemsep=0.1em,topsep=0.2em}
\usepackage{caption}
\captionsetup{font=small,labelfont=bf}
\usepackage{xcolor}
\usepackage{url}
\urlstyle{same}
\emergencystretch=3em
\sloppy

\newcommand{\Fq}{\mathbb{F}_q}
\newcommand{\Fp}{\mathbb{F}_p}
\newcommand{\Fr}{\mathbb{F}_r}
\newcommand{\Gas}{\mathrm{gas}}
\newcommand{\code}[1]{\texttt{\detokenize{#1}}}
\newcolumntype{Y}{>{\RaggedRight\arraybackslash}X}

\begin{document}

\begin{center}
    {\Large\bfseries Итоговый отчет о выполненной работе}\\[0.4em]
    {\large Исследование стоимости сопряжений на циклических кривых\\
    в EVM и перспектив дешевого folding}\\[0.4em]
    {\normalsize MNT4-753, MNT6-753 и исследовательская конструкция lollipop-305}
\end{center}

\tableofcontents
\vspace{0.4em}

\section{Постановка задачи и практическая мотивация}

Работа посвящена проверяемым вычислениям на эллиптических кривых в EVM.
Практическая мотивация связана с рекурсивными доказательствами. Обычный
SNARK подтверждает корректность одного вычисления. В рекурсивной схеме
следующее доказательство дополнительно проверяет предыдущее. Это позволяет
собирать длинную последовательность шагов в один компактный итоговый
результат. Такой подход нужен, например, для накопления состояний
распределенной системы, периодической публикации короткого доказательства
большого числа транзакций и построения проверяемых вычислений, которые
выполняются дольше одного блока.

Folding-схемы развивают ту же идею: вместо генерации независимого полного
доказательства после каждого шага последовательно объединяется отношение,
описывающее корректность вычислений. В конце проверяется компактный итог.
Поэтому стоимость одного шага folding критична: даже локально небольшой
избыточный расход умножается на длину всей цепочки.

В существующих реализациях, включая Sonobe/CycleFold, заметная часть стоимости
может возникать не из-за пользовательского утверждения, а из-за переноса
арифметики одной кривой в поле другой кривой. Такая арифметика называется
ненативной. Элемент большого поля раскладывается на машинные слова
(\emph{limb}-ы), а каждое умножение требует ограничений на частичные
произведения, переносы, проверку диапазонов и редукцию. В результате даже
простое пользовательское утверждение может сопровождаться тяжелой служебной
проверкой кривой. В качестве практического ориентира в работе используется
порядок около $9$ млн constraints для Sonobe/CycleFold-like decider.

Для Ethereum задача дополнительно осложняется архитектурой EVM:
\begin{itemize}
    \item EVM работает с 256-битными словами;
    \item для BN254 существуют предкомпилированные операции;
    \item для MNT4-753 и MNT6-753 аналогичных предкомпилированных операций нет;
    \item элементы MNT-полей имеют размер около 753 бит и требуют трех слов EVM.
\end{itemize}

Поэтому были исследованы два пути:
\begin{enumerate}
    \item выполнять арифметику сопряжения непосредственно on-chain и платить
    за нее gas;
    \item переносить вычисление off-chain и платить за доказательство,
    constraints, calldata и время работы.
\end{enumerate}

Потенциальный способ уменьшить число constraints состоит в том, чтобы
использовать цикл кривых: скалярное поле одной кривой совпадает с базовым полем
другой. Тогда часть арифметики следующего рекурсивного шага выражается
нативно, без дорогого переноса в постороннее поле. Однако отсутствие
MNT-precompile в EVM создает вторую сторону задачи: необходимо проверить,
можно ли реализовать такой путь с приемлемой on-chain стоимостью.

Цель работы состоит в том, чтобы реализовать и измерить основные варианты,
определить границу применимости и проверить, дает ли цикл MNT4/MNT6 или
конструкция lollipop более дешевую основу для будущего folding.

\section{Теоретическая база вычисления сопряжения}

\subsection{Цикл Миллера и финальная экспонента}

Для reduced Tate pairing используется вычисление функции Миллера и финальная
экспонента:
\[
    e(P,Q)=f_{r,P}(Q)^{(q^k-1)/r}.
\]
Для каждого шага обновляется аккумулятор:
\[
    f \leftarrow f^2\cdot \ell_{\mathrm{dbl}}(Q),
\]
а для ненулевой цифры signed-представления дополнительно:
\[
    f \leftarrow f\cdot \ell_{\mathrm{add}}(Q).
\]

Для уравнения сопряжений
\[
    e(P,Q)=e(R,S)
\]
удобно проверять эквивалентную форму
\[
    e(P,Q)\cdot e(-R,S)=1.
\]
Обе пары можно обрабатывать общим аккумулятором:
\[
    f \leftarrow
    f^2\cdot \ell_{Q,i}(P)\cdot \ell_{S,i}(-R).
\]
Главное преимущество общего аккумулятора: возведение в квадрат выполняется
один раз на раунд, а не отдельно для каждой пары (в некотором смысле паралельность операций).

\subsection{Удаление знаменателя}

В общей формуле функции Миллера возникает отношение линии к вертикальной
линии:
\[
    g_{A,B}(X)=\frac{\ell_{A,B}(X)}{v_{A+B}(X)}.
\]
Для MNT4-753 финальная экспонента убирает этот множитель. Поэтому в вычислительном пути
можно использовать только числитель линии. Практический смысл оптимизации существенен: на каждом шаге не требуется вычислять дополнительное обращение или вести второй аккумулятор знаменателя.
Это свойство не является автоматически верным для любой кривой. Например для третьей кривой lollipop-305 оно не выполняется в нужной форме, и это напрямую увеличивает стоимость.


\subsection{Проверка финальной экспоненты через свидетельство}

Статья \emph{On Proving Pairings}~\cite{eprint640} позволяет заменить полную
финальную экспоненту проверкой отношения. Для уравнения сопряжений пусть $F$ ---
результат объединенного цикла Миллера. Требуется проверить:
\[
    F^{(q^k-1)/r}=1.
\]
При выполнении условий теоремы статьи это равносильно существованию элемента
$c$, для которого:
\[
    F=c^r.
\]
Вместо длинной цепочки возведений в степень контракт проверяет эквивалентное
отношение:
\[
    F\cdot c^{-r}=1.
\]
Часть степени $c^{-r}$ встраивается в тот же signed double-and-add цикл, а
возведение в степень $q$ выполняется отображением Фробениуса.

\section{Оптимизированная 3-limb арифметика MNT4-753}

\subsection{Представление данных}

Элемент базового поля MNT4-753 не помещается в одно слово EVM и хранится как
три 256-битных слова:
\[
    x=x_0+2^{256}x_1+2^{512}x_2.
\]
Для длинных цепочек умножений используется Montgomery-представление.

\subsection{Почему выбран Montgomery/CIOS}

Были реализованы и измерены несколько вариантов:
\begin{itemize}
    \item Montgomery/CIOS;
    \item Barrett-редукция;
    \item FIOS;
    \item Comba/SOS;
    \item square-specialized Comba;
    \item branchless reduction;
    \item lazy reduction variants.
\end{itemize}

CIOS (\emph{Coarsely Integrated Operand Scanning}) объединяет накопление
частичных произведений и Montgomery-редукцию. Для EVM это выгоднее, чем
материализация полного 6-словного произведения в памяти.

\begin{table}[htbp]
\centering
\caption{Сравнение базовых 3-limb операций}
\small
\begin{tabularx}{\linewidth}{Y l r Y}
\toprule
Вариант & Операция & Gas/op & Вывод \\
\midrule
Montgomery/CIOS & $F_q$ mul & 2,959 & основной путь \\
Montgomery/CIOS & $F_q$ square & 2,947 & основной путь \\
Barrett & $F_q$ mul & 46,998 & существенно дороже \\
Barrett & $F_q$ square & 47,172 & существенно дороже \\
FIOS & $F_q$ mul & 18,509 & дороже CIOS \\
FIOS & $F_q$ square & 18,543 & дороже CIOS \\
Comba/SOS & $F_q$ mul & 18,301 & проигрывает из-за memory traffic \\
Square-Comba & $F_q$ square & 18,420 & лучше общего Comba, но хуже CIOS \\
\bottomrule
\end{tabularx}
\end{table}

\subsection{Оценка снизу и практическая оптимальность 3-limb умножения}

Для произведения двух трехсловных чисел требуется не менее девяти попарных
произведений:
\[
    ab=\sum_{i=0}^{2}\sum_{j=0}^{2}a_i b_j 2^{256(i+j)}.
\]
Montgomery-редукция выполняет три шага. На каждом шаге добавляется произведение
корректирующего коэффициента на трехсловный модуль, то есть еще три произведения
на шаг. Получаем:
\[
    9\text{ произведений для }ab+
    9\text{ произведений для REDC}=18
\]
полных 512-битных limb-произведений.

В EVM нет инструкции \code{MULHI}. Для восстановления старшей половины
512-битного произведения используется связка:
\begin{verbatim}
lo := mul(u, v)
mm := mulmod(u, v, not(0))
hi := sub(sub(mm, lo), lt(mm, lo))
\end{verbatim}

При действующем расписании EVM опкоды имеют стоимость:
\[
    \operatorname{MUL}=5,\quad
    \operatorname{MULMOD}=8,\quad
    \operatorname{ADD}=\operatorname{SUB}=\operatorname{LT}=3.
\]
Следовательно, одно восстановление полного 512-битного произведения требует
как минимум:
\[
    5+8+3+3+3=22\ \Gas.
\]
Чтобы оценка была содержательной, в работе дополнительно построена  opcode-трасса оптимизированного all-stack Yul-кода.

Абстрактная схема CIOS содержит девять limb-произведений $a_i b_j$ и девять
произведений $m_i q_j$ при редукции. Однако младшие произведения
$m_i q_0$ не требуется вычислять полностью: после выбора
\[
    m_i=t_0\cdot(-q^{-1})\bmod 2^{256}
\]
младшее слово обнуляется по модулю $2^{256}$ и затем отбрасывается при сдвиге
окна CIOS. В реализации сохранено только вычисление старшей половины этих
трех произведений. Поэтому фактический оптимизированный граф содержит:
\[
\begin{aligned}
    &18 &&\text{инструкций \code{MULMOD} для старших половин;}\\
    &15 &&\text{полезных младших limb-произведений \code{MUL};}\\
    &3  &&\text{инструкции \code{MUL} для коэффициентов }m_i.
\end{aligned}
\]

Для воспроизводимой проверки выполнен одиночный внешний вызов
\code{montMul3} с full-width входами и снята geth-style opcode-трасса. Она
содержит $1{,}012$ исполняемых инструкций и расходует $3{,}316$ gas без
учета базовой стоимости транзакции. Из них $2{,}843$ gas относятся к
развернутому CIOS-телу функции; оставшиеся $473$ gas приходятся на диспетчер
внешнего вызова, ABI-декодирование и возврат результата.

Структура исполняемого CIOS-тела приведена в
табл.~\ref{tab:montmul3-opcode-structure}.

\begin{table}[htbp]
\centering
\caption{Динамический opcode-состав развернутого CIOS-тела \code{montMul3}}
\label{tab:montmul3-opcode-structure}
\small
\begin{tabular}{l r r}
\toprule
Группа & Число инструкций & Сумма, gas \\
\midrule
\code{MUL} & 18 & 90 \\
\code{MULMOD} & 18 & 144 \\
\code{ADD} & 63 & 189 \\
\code{SUB} & 36 & 108 \\
\code{LT} & 61 & 183 \\
\code{JUMPI} & 19 & 190 \\
\code{MLOAD}, \code{MSTORE}, \code{CODECOPY} & $40+20+20$ & 300 \\
\code{PUSH}, \code{DUP}, \code{SWAP} и остальные & --- & 1,639 \\
\midrule
Итого: развернутое CIOS-тело & --- & 2,843 \\
\bottomrule
\end{tabular}
\end{table}

Выбранный вариант является лучшим подтвержденным практическим путем в исследованном классе алгоритмов, а неиспользованного алгоритмического резерва порядка десятков процентов не обнаружено.

\subsection{Башня расширений и Karatsuba}

Для MNT4 используется башня:
\[
    \mathbb{F}_{q^2}=\Fq[u]/(u^2-13),
    \qquad
    \mathbb{F}_{q^4}=\mathbb{F}_{q^2}[v]/(v^2-u).
\]
Умножения на фиксированные элементы поля, не являющиеся квадратами в соответствующих расширениях, реализованы с помощью специализированных формул. Например:
\[
    (x_0+x_1u)u=13x_1+x_0u.
\]

\begin{table}[htbp]
\centering
\caption{Стоимость операций расширений MNT4-753}
\small
\begin{tabular}{l r}
\toprule
Операция & Gas/op \\
\midrule
$\Fq$ mul & 2,959 \\
$\Fq$ square & 2,947 \\
$\mathbb{F}_{q^2}$ mul & 11,877 \\
$\mathbb{F}_{q^2}$ square & 10,592 \\
$\mathbb{F}_{q^4}$ mul & 42,764 \\
$\mathbb{F}_{q^4}$ square & 36,606 \\
\bottomrule
\end{tabular}
\end{table}

\begin{table}[htbp]
\centering
\caption{Эффект специализированных умножений}
\small
\begin{tabular}{l r r r}
\toprule
Операция & Общий путь & Специализированный путь & Выигрыш \\
\midrule
умножение на $13$ в $\Fq$ & 2,961 & 1,481 & $2.0\times$ \\
умножение на $u$ в $\mathbb{F}_{q^2}$ & 12,587 & 1,731 & $7.3\times$ \\
умножение на $v$ в $\mathbb{F}_{q^4}$ & 42,253 & 1,859 & $22.7\times$ \\
\bottomrule
\end{tabular}
\end{table}

\subsection{Подготовленные разреженные линии}

После подстановки точки линия Миллера имеет специальную структуру. Ее не нужно
представлять как произвольный плотный элемент $\mathbb{F}_{q^4}$. В проекте
реализованы:
\begin{itemize}
    \item prepared fixed-$Q$ cache;
    \item sparse-формат коэффициентов;
    \item специализированное умножение аккумулятора на линию;
    \item pointer/packed memory API;
    \item потоковое чтение code-shards через \code{EXTCODECOPY};
    \item экспериментальное слияние $f\leftarrow f^2\ell(P)$ без лишнего
    промежуточного объекта.
\end{itemize}

\paragraph{Формальная идея разреженного умножения.}
Пусть аккумулятор цикла Миллера имеет общий вид:
\[
    f=f_0+f_1v,\qquad f_0,f_1\in\mathbb{F}_{q^2}.
\]
После подстановки $G_1$-точки подготовленная линия не является произвольным
элементом $\mathbb{F}_{q^4}$: часть ее коэффициентов равна нулю или выражается
через заранее известные координаты. Условно ее можно записать как:
\[
    \ell(P)=\ell_0(P)+\ell_1(P)v,
\]
где $\ell_0,\ell_1$ вычисляются из компактного набора prepared-коэффициентов.
Вместо построения плотного объекта и вызова общего умножения
\[
    f\leftarrow\operatorname{Mul}_{\mathbb{F}_{q^4}}(f,\ell(P))
\]
используется специализированная формула:
\[
\begin{aligned}
    t_0&=f_0\ell_0,\\
    t_1&=f_1\ell_1,\\
    t_2&=(f_0+f_1)(\ell_0+\ell_1),\\
    f'_0&=t_0+u\,t_1,\\
    f'_1&=t_2-t_0-t_1.
\end{aligned}
\]
Умножение на $u$ выполняется дешевым переставлением коэффициентов. 

\subsection{Низкоуровневые оптимизации EVM}

Алгоритмические формулы были дополнены низкоуровневой реализацией:
\begin{itemize}
    \item фиксированный размер $3\times3$ limb развернут вручную, без
    динамических циклов и динамических массивов;
    \item старшая половина произведения восстанавливается через
    \code{MUL}/\code{MULMOD}
    \item CIOS объединяет умножение и REDC, чтобы не сохранять полное
    шестисловное произведение;
    \item промежуточные значения по возможности удерживаются на стеке;
    \item pointer/packed API и scratch arena уменьшают копирование вложенных
    структур расширений;
    \item lazy reduction применяется только там, где диапазон промежуточного
    значения формально ограничен и отложенная редукция действительно дешевле;
    \item специальные операции \code{mulBy13}, \code{mulByU},
    \code{mulByV} заменяют общие умножения;
    \item prepared-кэш может читаться потоково из data-контрактов через
    \code{EXTCODECOPY}, без повторной передачи большого blob;
\end{itemize}

\section{MNT4: от наивного Tate pairing к оптимизированному пути}

\subsection{Начальная точка отсчета}

Для честного сравнения построена наивная модель Tate pairing. Она намеренно не
использует:
\begin{itemize}
    \item короткий Ate loop;
    \item prepared lines;
    \item sparse multiplication;
    \item Yul hot path;
    \item Frobenius-декомпозицию финальной экспоненты.
\end{itemize}

Для такого вызова построена нижняя экстраполяция:
\[
\begin{aligned}
    \text{Miller} &=
    752\cdot 973{,}261+372\cdot 486{,}819\\
    &=912{,}988{,}940\ \Gas,\\
    \text{FE} &=
    2259\cdot 486{,}899+1100\cdot 486{,}819\\
    &=1{,}635{,}405{,}741\ \Gas.
\end{aligned}
\]
Итого:
\[
    2{,}548{,}394{,}681\ \Gas.
\]
Это нижняя оценка: она не включает построение линий, проверки входов и часть
накладных расходов.

\subsection{Оптимизационная лестница}

Не все оптимизации можно честно представить как последовательность полных
контрактов: некоторые меняют стоимость одной операции поля, другие ---
длину цикла Миллера, третьи --- способ доставки кэша. Поэтому ниже приведены
три взаимосвязанные таблицы.

\begin{table}[htbp]
\centering
\caption{Лестница полного MNT4-пути}
\small
\begin{tabularx}{\linewidth}{r Y r r}
\toprule
Шаг & Реализация & Gas & Снижение к предыдущему шагу \\
\midrule
0 & Наивная Tate cost model & 2,548,394,681 & --- \\
1 & Полный optimized fixed-$Q$ on-chain путь & 258,753,182 & 89.85\% \\
2 & Fixed-$Q$ prepared sparse blob & 79,588,799 & 69.24\% \\
3 & Fixed-$Q$ prepared sparse code-shards & 79,586,596 & 0.003\% \\
\bottomrule
\end{tabularx}
\end{table}

\begin{table}[htbp]
\centering
\caption{Локальные оптимизации арифметики}
\small
\begin{tabularx}{\linewidth}{Y Y r r}
\toprule
Оптимизация & Что меняется & Было & Стало \\
\midrule
Montgomery/CIOS + Yul & $F_q$ mul вместо generic Solidity multi-limb &
18,770 & 2,959 \\
Karatsuba для башни & $F_{q^4}$ mul вместо generic расширения &
486,819 & 42,764 \\
Умножение на $13$ & общая операция заменяется сложениями &
2,961 & 1,481 \\
Умножение на $u$ & $(x_0+x_1u)u=13x_1+x_0u$ &
12,587 & 1,731 \\
Умножение на $v$ & перестановка $F_{q^2}$-коэффициентов &
42,253 & 1,859 \\
\bottomrule
\end{tabularx}
\end{table}

\begin{table}[htbp]
\centering
\caption{Лестница цикла Миллера, финальной экспоненты и кэша}
\small
\begin{tabularx}{\linewidth}{Y Y r r}
\toprule
Оптимизация & Смысл & Было & Стало \\
\midrule
Ate вместо полного Tate loop & короткая signed-запись параметра &
912,988,940 & 427,284,953 \\
Sparse line multiplication & линия не материализуется как общий $F_{q^4}$ &
973,261/шаг & $\approx137,874$/digit \\
Prepared fixed-$Q$ cache & line generation переносится из on-chain &
258,753,182 & 79,588,799 \\
Frobenius/$w_0$ FE & декомпозиция вместо binary exponentiation &
1,635,405,741 & $\approx28,134,825$ \\
Aggressive fused hot path & слияние $f^2$ и sparse line multiplication &
51,978,344 & 51,949,399 \\
Code-shards & повторный blob заменяется \code{EXTCODECOPY} &
79,588,799 & 79,586,596 \\
\bottomrule
\end{tabularx}
\end{table}

\subsection{Формальная нижняя оценка стоимости Miller core}

Для двухпарного уравнения:
\[
    e(P,Q)e(-R,S)=1
\]
используем:
\[
    L=376,\qquad H=123,
\]
где $L$ --- число удвоений, $H$ --- число ненулевых addition-шагов.
Минимально нужны:
\begin{itemize}
    \item $L$ возведений аккумулятора в квадрат;
    \item $2L$ применений doubling-линий;
    \item $2H$ применений addition-линий.
\end{itemize}
Грубая нижняя оценка через измеренные $\mathbb{F}_{q^4}$-операции:
\[
\begin{aligned}
    376\cdot 36{,}606 &= 13{,}763{,}856,\\
    2\cdot376\cdot42{,}764 &= 32{,}157{,}728,\\
    2\cdot123\cdot42{,}764 &= 10{,}499{,}864.
\end{aligned}
\]
Суммарно:
\[
    56{,}421{,}448\ \Gas.
\]
Реальный код дополнительно загружает sparse-коэффициенты, вычисляет значения
линий, работает с памятью и обрабатывает signed loop. Поэтому результат порядка $87.7$ млн gas согласуется с формальной моделью.

\section{MNT4 direct residue verifier по ePrint 2024/640}

\subsection{API и граница доверия}

В основном fixed-shards режиме контракт проверяет:
\[
    e(P,Q)e(-R,S)=1,
\]
где $Q,S\in G_2$ фиксированы при развертывании, а $P,R\in G_1$ передаются
пользователем. Prepared-кэш линий хранится в runtime-коде data-контрактов.
Адреса shards фиксируются в конструкторе. Пользователь не может заменить их
при вызове.

\subsection{Результаты MNT4}

\begin{table}[htbp]
\centering
\caption{Direct residue verifier MNT4}
\small
\begin{tabularx}{\linewidth}{Y Y r}
\toprule
Сценарий & Что выполняется & Gas \\
\midrule
Miller core & цикл Миллера без финальной экспоненты & 87,747,219 \\
Full equation & Miller core + обычная FE & 115,997,313 \\
Residue equation & Miller core + $c$-проверка вместо полной FE & 93,879,746 \\
\bottomrule
\end{tabularx}
\end{table}

Вклад обычной FE:
\[
    115{,}997{,}313-87{,}747{,}219
    =28{,}250{,}094\ \Gas.
\]
Вклад residue-проверки:
\[
    93{,}879{,}746-87{,}747{,}219
    =6{,}132{,}527\ \Gas.
\]
Экономия:
\[
    22{,}117{,}567\ \Gas.
\]
Оптимизация финальной экспоненты работает, но не устраняет главный расход:
цикл Миллера.

\section{Проверка цикла Миллера через полиномиальные отношения}

\subsection{Математическая идея}

Следующий шаг статьи~\cite{eprint640} --- не исполнять все умножения
расширенного поля on-chain, а проверить их как полиномиальные отношения.
В статье сформулирована общая идея. Для MNT4-753 потребовалась отдельная
адаптация: нужно связать quotient-check не с абстрактными умножениями, а с
реальными состояниями общего multi-Miller аккумулятора, fixed-$Q$/fixed-$S$
кэшем и residue-свидетельством $c$.

Пусть:
\[
    \mathbb{F}_{q^4}\simeq \Fq[X]/(p(X)).
\]
Для умножения $a\cdot b=c$ существует частный многочлен $Q(X)$:
\[
    a(X)b(X)-c(X)=Q(X)p(X).
\]
Verifier выбирает случайную точку $\alpha$ и проверяет:
\[
    a(\alpha)b(\alpha)-c(\alpha)=Q(\alpha)p(\alpha).
\]
Для набора отношений используется challenge $\beta$:
\[
    \sum_i\beta^i R_i(X)=Q(X)Z_H(X).
\]
Для цикла Миллера $R_i$ кодируют реальные переходы:
\[
    f_{i+1}=f_i^2\ell_{Q,i}(P)\ell_{S,i}(-R).
\]

Чтобы prover не подменил значения после получения $\alpha$, нужны commitments
и доказательства открытия.

\subsection{Построенная MNT4-модель}

Для MNT4 была построена следующая модель.

\paragraph{Блочное сжатие цикла.}
Несколько раундов объединяются в один блок. Для двух шагов:
\[
    f_{i+2}
    =
    f_i^4\ell_i^2\ell_{i+1}.
\]
Для блока длины $d$:
\[
    f_{i+d}
    =
    f_i^{2^d}
    \prod_{j=0}^{d-1}\ell_{i+j}^{2^{d-1-j}}.
\]
В модели используется $d=5$. Тогда один блочный переход для уравнения
сопряжений имеет вид:
\[
    F_{b+1}
    =
    F_b^{32}
    G_b^Q(P)
    G_b^S(-R)
    c^{\kappa_b}.
\]
Здесь $G_b^Q$ и $G_b^S$ --- сжатые делители для зафиксированных точек
$Q,S$, а $\kappa_b$ --- заранее известный показатель residue-рекурсии.

\paragraph{Проверка вычисления делителей.}
Для fixed-$Q$/fixed-$S$ коэффициенты делителей строятся off-chain и
фиксируются root-ом. Значение $G_b^Q(P)$ не должно выбираться prover-ом
произвольно. Поэтому строится таблица вычисления по схеме Горнера:
\[
    h_{j+1}=h_j P_x+\gamma_j,
\]
где $\gamma_j$ --- зафиксированный коэффициент сжатого делителя. Аналогично
вычисляется $G_b^S(-R)$. Локальные отношения Горнера и блочные переходы
собираются в quotient polynomial.

\paragraph{Что доказано для модели.}
Для спецификации были формализованы:
\begin{itemize}
    \item связь пяти пошаговых переходов Миллера с одним сжатым делителем;
    \item корректность Horner-рекурсии для fixed-divisor коэффициентов;
    \item quotient identity для локальных переходов;
    \item Merkle-binding динамических таблиц;
    \item Fiat--Shamir transcript и детерминированный порядок запросов;
    \item ordinary-FRI low-degree проверка;
    \item итоговая soundness-оценка как сумма ошибок FRI,
    Schwartz--Zippel, Merkle-binding и Fiat--Shamir модели.
\end{itemize}

\subsection{KZG over BN254}

KZG дает короткое доказательство открытия и дешевую on-chain проверку благодаря
BN254 precompile. В проекте реализован реальный opening layer:
\[
    C_j=\operatorname{Commit}(w_j),\qquad
    C_Q=\operatorname{Commit}(Q).
\]
Для одного полинома проверяется стандартное отношение:
\[
    e(C-yG_1+x\pi,G_2)=e(\pi,\tau G_2).
\]
Здесь $C$ --- commitment, $x$ --- точка открытия, $y=f(x)$ ---
заявленное значение, $\pi$ --- opening proof, $\tau G_2$ --- элемент SRS.
Measured стоимость одной on-chain opening-проверки составляет около:
\[
    133{,}039\ \Gas.
\]

Проблема переносится внутрь circuit: MNT4-арифметика становится non-native
относительно BN254.

\begin{table}[htbp]
\centering
\caption{Оценка BN254 non-native constraints}
\small
\begin{tabular}{l r}
\toprule
Операция & Constraints \\
\midrule
MNT4 $\Fq$ mul как non-native & 3,692 \\
MNT4 $\mathbb{F}_{q^2}$ mul как non-native & 11,300 \\
Модель MNT4 $\mathbb{F}_{q^4}$ mul как non-native & 63,436 \\
Приближенная sparse Miller relation & 63,309,128 \\
\bottomrule
\end{tabular}
\end{table}

Число $63{,}309{,}128$ получается не из размера KZG proof. Оно относится к
relation, которую нужно доказать до открытия: каждое MNT4 $F_{q^4}$
умножение при переносе в BN254 R1CS раскладывается на non-native операции.
В проекте отдельно измерены $3{,}692$ constraints для одного MNT4
$F_q$-умножения через emulated field, $11{,}300$ для $F_{q^2}$ и
$63{,}436$ для модели $F_{q^4}$. Подстановка числа sparse Miller
переходов дает ориентир $63.3$ млн constraints. То есть KZG уменьшает
gas verifier-а, но возвращает исходную проблему роста off-chain constraints.

\subsection{Merkle/FRI}

В альтернативном направлении prover коммитится к таблицам через Merkle root.
Verifier проверяет Merkle-пути открытых строк. FRI
(\emph{Fast Reed--Solomon Interactive Oracle Proof of Proximity}) доказывает,
что открываемые таблицы согласуются с многочленами ограниченной степени.
Fiat--Shamir-преобразование делает схему неинтерактивной: challenges
детерминированно вычисляются из transcript.

В проекте реализованы:
\begin{itemize}
    \item Merkle opening layer;
    \item исполняемый FRI эксперимент;
    \item спецификация блочно-сжатого ordinary-FRI пути;
    \item FRI cost model для нативного MNT4-поля.
\end{itemize}

Для MNT4 один элемент поля занимает:
\[
    3\cdot 32=96\text{ байт}.
\]
Умеренная модель открытий:
\[
    4096\cdot96+128\cdot16\cdot32+16\cdot32
    =459{,}264\text{ байт}.
\]
Актуальная FRI модель для 128-битного профиля:

\begin{table}[htbp]
\centering
\caption{Ordinary-FRI cost model}
\small
\begin{tabular}{l r}
\toprule
Показатель & Значение \\
\midrule
Число запросов & 576 \\
Оценка soundness & 135.76 бит \\
FRI schedule & $[2,2,2,4]$ \\
Размер доказательства & 2,072,224 байт \\
Нижняя оценка & 51,359,352 gas \\
Ожидаемая модельная оценка & 78,340,624 gas \\
\bottomrule
\end{tabular}
\end{table}

\paragraph{Из чего складывается нижняя оценка.}
Для каждого случайного запроса verifier должен получить раскрытые значения
столбцов, соседние значения для локального перехода, Merkle-path и элементы
слоев FRI. Поскольку один MNT4-элемент занимает $96$ байт, даже до исполнения
арифметики calldata становится существенным компонентом стоимости. Актуальная
модель суммирует:
\[
    \Gas_{\min}
    =
    \Gas_{\mathrm{calldata}}
    +
    \Gas_{\mathrm{Merkle}}
    +
    \Gas_{\mathrm{FRI\ folding}}
    +
    \Gas_{\mathrm{local\ checks}}.
\]
Для профиля с $576$ запросами строгая нижняя оценка равна
$51{,}359{,}352$ gas. Исполняемая модель с реалистичными накладными расходами дает:
\[
    78{,}340{,}624\ \Gas.
\]

\section{MNT6-753 и цикл MNT4/MNT6}

\subsection{MNT6 arithmetic и residue verifier}

Для MNT6 реализована башня:
\[
    \Fq\longrightarrow\mathbb{F}_{q^3}\longrightarrow\mathbb{F}_{q^6}.
\]
Реализованы packed/pointer API, scratch arena, prepared cache, потоковое чтение
code-shards, fused line multiplication и Frobenius/w0 FE.

\begin{table}[htbp]
\centering
\caption{Базовые операции MNT6}
\small
\begin{tabular}{l r}
\toprule
Операция & Gas/op \\
\midrule
$\Fq$ mul & 2,314 \\
$\Fq$ square & 2,301 \\
$\mathbb{F}_{q^3}$ mul & 36,472 \\
$\mathbb{F}_{q^3}$ square & 26,813 \\
$\mathbb{F}_{q^6}$ mul & 166,774 \\
$\mathbb{F}_{q^6}$ square & 113,788 \\
\bottomrule
\end{tabular}
\end{table}

\paragraph{Лемма 2. Общий аккумулятор для уравнения сопряжений.}
Пусть на раунде $i$ индивидуальные аккумуляторы обновляются как:
\[
    f'_{Q}=f_Q^2\ell_{Q,i}(P),
    \qquad
    f'_{S}=f_S^2\ell_{S,i}(-R).
\]
Для произведения $F=f_Qf_S$ имеем:
\[
    F'
    =
    f'_Qf'_S
    =
    (f_Qf_S)^2
    \ell_{Q,i}(P)\ell_{S,i}(-R).
\]
Следовательно, вместо двух возведений аккумуляторов в квадрат достаточно
одного.

\begin{table}[htbp]
\centering
\caption{Результаты MNT6}
\small
\begin{tabularx}{\linewidth}{Y r}
\toprule
Режим & Gas \\
\midrule
Miller loop, packed pointer blob & 93,254,054 \\
Packed Frobenius/w0 FE & 38,428,108 \\
Полное single pairing: Miller + packed FE & 131,685,843 \\
Диагностический residue digest одного pairing & 103,277,505 \\
Два Miller loop + полная FE для equation & 226,078,963 \\
Общий multi-Miller + $c$-проверка для equation & 172,004,717 \\
\bottomrule
\end{tabularx}
\end{table}

Экономия общего residue verifier относительно контрольного equation baseline:
\[
    226{,}078{,}963-172{,}004{,}717
    =54{,}074{,}246\ \Gas,
\]
то есть $23.92\%$.

\subsection{Cycle-native accounting}

Rust-модуль проверяет параметры цикла через arkworks, строит reference
pairings и считает operation/constraints accounting для будущего native
relation layer.


\begin{table}[htbp]
\centering
\caption{Предварительная модель MNT-cycle-native relation}
\small
\begin{tabular}{l l r r r}
\toprule
Сторона & Башня & Miller rounds & Additions & Prepared relation \\
\midrule
MNT4 & $\mathbb{F}_{q^2}/\mathbb{F}_{q^4}$ & 376 & 124 & 24,126 \\
MNT6 & $\mathbb{F}_{q^3}/\mathbb{F}_{q^6}$ & 376 & 123 & 48,942 \\
\midrule
\multicolumn{4}{r}{Сумма} & 73,068 \\
\bottomrule
\end{tabular}
\end{table}

Для сравнения:
\[
    \text{BN254 non-native sparse Miller estimate}=63{,}309{,}128,
\]
\[
    \text{Sonobe/CycleFold-like anchor}\approx9{,}000{,}000.
\]


\paragraph{Оценка будущего cycle-native folding слоя.}
Внутреннее ядро relation уже посчитано по формулам проекта:
\[
    24{,}126+48{,}942=73{,}068
\]
нативных multiplication constraints. 

\begin{table}[htbp]
\centering
\caption{Проекция cycle-native constraints при разных накладных расходах}
\small
\begin{tabular}{l r r}
\toprule
Сценарий & Коэффициент к relation core & Constraints \\
\midrule
Только подготовленные fragments & $1\times$ & 73,068 \\
Легкая folding-обвязка & $2\times$ & 146,136 \\
Умеренная folding-обвязка & $4\times$ & 292,272 \\
Консервативная folding-обвязка & $8\times$ & 584,544 \\
Верхняя исследовательская граница & $12\times$ & 876,816 \\
\bottomrule
\end{tabular}
\end{table}

\section{Исследовательское направление lollipop-305}

\subsection{Идея}

Статья \emph{Cycles of supersingular elliptic curves for pairing-based proof
systems}~\cite{eprint1627} предлагает конструкции, в которых часть
арифметики выполняется над меньшими полями. Для исследовательского примера
lollipop-305 выбран параметр:
\[
    x=8004046504391788107635887004283725454478544674,
\]
\[
    p=x^2-x+1,\qquad q=x^2+1.
\]
Оба основных поля имеют размер около 305 бит и помещаются в два слова EVM.
Цель этого направления --- проверить, является ли уменьшение числа limb-ов
достаточным условием практического выигрыша. 

\subsection{Кривые конструкции}

\begin{table}[htbp]
\centering
\caption{Pairing-компоненты lollipop-305}
\small
\begin{tabularx}{\linewidth}{l Y Y}
\toprule
Часть & Кривая и поле & Целевое поле \\
\midrule
stick & pairing-friendly prototype над $\Fp$ & $\mathbb{F}_{p^4}$ \\
cycle $E$ & $E/\mathbb{F}_{p^2}: y^2=x^3+(\mu+1)x$ & $\mathbb{F}_{p^4}$ \\
cycle $\widehat{E}$ &
$\widehat{E}/\mathbb{F}_{q^2}: y^2=x^3+(\eta+2)$ &
$\mathbb{F}_{q^6}$ \\
\bottomrule
\end{tabularx}
\end{table}

Полная конструкция содержит три pairing-компонента:
\begin{enumerate}
    \item ``палочку'' --- pairing-friendly prototype над $\Fp$;
    \item первую supersingular cycle-кривую $E/\mathbb{F}_{p^2}$;
    \item вторую supersingular cycle-кривую
    $\widehat{E}/\mathbb{F}_{q^2}$.
\end{enumerate}
Дополнительно статья использует меньшую непаринговую кривую для рекурсивного
сценария, но она не является отдельным EVM pairing verifier.

Проверены отношения:
\[
    \#E(\mathbb{F}_{p^2})=p^2+1=qN_q,
\]
\[
    \#\widehat{E}(\mathbb{F}_{q^2})=q^2-q+1=pN_p.
\]

\paragraph{Лемма 3. Исправление башни для stick/$E$-части.}
В работе вычислительно найден и проверен неквадрат:
\[
    \xi=1+u.
\]
Поэтому используется башня:
\[
    \mathbb{F}_{p^4}
    =
    \mathbb{F}_{p^2}[v]/(v^2-\xi).
\]

\paragraph{Лемма 4. Twist/Frobenius-отношение.}
Для башни построено отображение из twist в исходную кривую:
\[
    \psi(X,Y)=(\xi X,\xi vY).
\]
backend проверяет:
\[
    [r]P=\mathcal O,\quad
    [r]\psi(Q')=\mathcal O,\quad
    \pi_p(\psi(Q'))=[p]\psi(Q'),
\]
а также нетривиальность результата сопряжения и условие принадлежности
целевой подгруппе.

\paragraph{Лемма 5. Корректная форма для $\widehat E$.}
Для третьей кривой выбрана башня:
\[
    \mathbb{F}_{q^6}
    =
    \mathbb{F}_{q^2}[w]/(w^3-\rho),
\]
где:
\[
    \rho^2=\frac{2+\eta}{2-\eta}.
\]

Для $\widehat E$ упрощенная Tate-style формула не проходит. Корректным
reference путем является Weil-отношение:
\[
    e_p(P,Q)
    =
    (-1)^p\frac{f_{p,P}(Q)}{f_{p,Q}(P)}.
\]
Отсюда следует verifier relation:
\[
    f_{p,P}(Q)f_{p,-P}(Q)
    =
    f_{p,Q}(P)f_{p,Q}(-P).
\]

\subsection{Оптимизированная 2-limb арифметика}

Для двухсловного Montgomery-умножения generic lower bound содержит:
\[
    4\text{ произведения для }ab+
    4\text{ произведения для REDC}.
\]
Для lollipop-305 старший limb модуля меньше $2^{49}$. Поэтому:
\[
    a_1b_1<2^{98}<2^{256}.
\]
Одно полное 512-битное восстановление можно заменить обычным \code{MUL}.
Дополнительно убрана запись limb-а, который сразу удаляется CIOS-сдвигом.

Для 305-битных полей проделан тот же класс работы, что и для 753-битной арифметики MNT4:
\begin{itemize}
    \item Montgomery/CIOS-представление;
    \item ручной Yul hot path;
    \item специальная обработка малого старшего limb-а;
    \item арифметика расширений $\mathbb{F}_{p^2}$,
    $\mathbb{F}_{p^4}$, $\mathbb{F}_{q^2}$ и $\mathbb{F}_{q^6}$;
    \item Karatsuba и специализированные умножения на нерезиденты;
    \item prepared lines, residue-пути и fixed-cache режимы;
\end{itemize}

Для наиболее дорогой части $\widehat E/\mathbb{F}_{q^2}$ дополнительно
выполнен отдельный hot-path аудит. Для
\[
    a=a_0+a_1w+a_2w^2,\qquad w^3=\rho
\]
общий вызов умножения $a\cdot a$ заменен специализированным квадратом:
\[
\begin{aligned}
    c_0&=a_0^2+2\rho a_1a_2,\\
    c_1&=2a_0a_1+\rho a_2^2,\\
    c_2&=a_1^2+2a_0a_2.
\end{aligned}
\]
Локальная стоимость квадрата $\mathbb{F}_{q^6}$ снизилась с $43{,}972$ до
$37{,}373$ gas/op. В Solidity verifier также перенесены единая scratch arena
для временных значений и перестановка указателей аккумуляторов вместо полного
копирования после каждого шага. Потоковое чтение code-shards по одной линии и
pointer-вариант для $\mathbb{F}_{p^4}$ были реализованы для сравнения, но
оказались дороже и не включены в основной путь.

\begin{table}[htbp]
\centering
\caption{Сравнение 2-limb lollipop-305 и 3-limb MNT4}
\small
\begin{tabular}{l r r r}
\toprule
Операция & lollipop-305 & MNT4-753 & Выигрыш \\
\midrule
$F_p$ mul & 1,018 & 2,959 & $2.91\times$ \\
$F_p$ square & 940 & 2,947 & $3.14\times$ \\
$F_{p^2}$ mul & 4,171 & 11,877 & $2.85\times$ \\
$F_{p^2}$ square & 3,254 & 10,592 & $3.26\times$ \\
$F_{p^4}$ mul & 14,942 & 42,764 & $2.86\times$ \\
$F_{p^4}$ square & 13,211 & 36,606 & $2.77\times$ \\
\bottomrule
\end{tabular}
\end{table}

\subsection{Почему для третьей кривой нужны два аккумулятора}

Для stick и $E$ можно использовать Article640-style residue verifier. Для
$\widehat{E}/\mathbb{F}_{q^2}$ степень вложения равна $3$. Стандартное удаление
знаменателей линий не работает в требуемой форме. Поэтому backend и Solidity
отдельно накапливают:
\[
    F_{\mathrm{num}},\qquad F_{\mathrm{den}}.
\]
На doubling-шаге:
\[
    F_{\mathrm{num}}\leftarrow
    F_{\mathrm{num}}^2\ell_{\mathrm{dbl}}(P),
\]
\[
    F_{\mathrm{den}}\leftarrow
    F_{\mathrm{den}}^2v_{\mathrm{dbl}}(P).
\]
На addition-шаге:
\[
    F_{\mathrm{num}}\leftarrow
    F_{\mathrm{num}}\ell_{\mathrm{add}}(P),
\]
\[
    F_{\mathrm{den}}\leftarrow
    F_{\mathrm{den}}v_{\mathrm{add}}(P).
\]
В конце проверяется отношение:
\[
    c^p F_{\mathrm{den}}=F_{\mathrm{num}}.
\]
Это эквивалентно residue-проверке частного:
\[
    F=F_{\mathrm{num}}/F_{\mathrm{den}}.
\]
Ведение числителя и знаменателя остается необходимым, но итоговая реализация
не запускает две независимые трассы для $P$ и $-P$. Вместо этого используется
общий product-accumulator: числитель умножается на линии для $P$ и $-P$, а
знаменатель, зависящий только от $x(P)$, ведется один раз и в конце
возводится в квадрат. Дополнительно длинное возведение $c^p$ заменено
декомпозицией
\[
    c^p=c^{q-(q-p)}=c^q\cdot(c^{-1})^{q-p}.
\]
Значение $c^q$ вычисляется отображением Фробениуса в $\mathbb{F}_{q^6}$, а
оставшаяся степень $q-p$ имеет длину только $153$ бита. Контракт получает
$c^{-1}$ как witness и проверяет $c\cdot c^{-1}=1$.

\subsection{Результаты lollipop-305}

\begin{table}[htbp]
\centering
\caption{Pairing-verifier-ы lollipop-305}
\small
\begin{tabularx}{\linewidth}{Y Y r}
\toprule
Часть & Режим & Function gas \\
\midrule
stick & residue FE & 8,700,515 \\
cycle $E/\mathbb{F}_{p^2}$ & residue FE & 18,340,638 \\
cycle $\widehat{E}/\mathbb{F}_{q^2}$ & Weil fallback & 174,865,299 \\
cycle $\widehat{E}/\mathbb{F}_{q^2}$ & raw product trace $F_{\mathrm{num}},F_{\mathrm{den}}$ & 44,734,048 \\
cycle $\widehat{E}/\mathbb{F}_{q^2}$ & product trace + Frobenius residue & 56,163,308 \\
\midrule
Полный lollipop-cycle & stick + $E$ + оптимизированный $\widehat{E}$ & 83,270,967 \\
\bottomrule
\end{tabularx}
\end{table}

Для fixed-cache режима основным значением является полный verifier-вызов
\code{verifyEhatAteResidueProductFrobeniusFixedShards}: он читает кэш из
code-shards, строит $F_{\mathrm{num}}$ и $F_{\mathrm{den}}$, проверяет
$c\cdot c^{-1}=1$, вычисляет $c^p$ через Фробениус-декомпозицию и проверяет
$c^pF_{\mathrm{den}}=F_{\mathrm{num}}$. Его стоимость равна
\[
    56{,}163{,}308\ \Gas.
\]
Отдельное значение $44{,}734{,}048$ gas относится только к raw product trace:
это вклад оптимизированного цикла Миллера без финальной проверки с
$c$-свидетельством и без fixed-shards wrapper. Поэтому для итоговых таблиц
используется $56{,}163{,}308$ gas, а $44{,}734{,}048$ gas указывается только
как разбиение стоимости.

\section{Сводные результаты и trade-off}

\begin{longtable}{p{0.31\linewidth}p{0.36\linewidth}r}
\caption{Основные результаты работы}\\
\toprule
Подход & Что измеряется & Результат \\
\midrule
\endfirsthead
\toprule
Подход & Что измеряется & Результат \\
\midrule
\endhead
Наивный Tate baseline & строгая нижняя экстраполяция & 2.55 млрд gas \\
MNT4 full optimized on-chain & fixed-$Q$, линии строятся on-chain & 258.8 млн gas \\
MNT4 Article640 residue & fixed-shards equation с проверкой $G_1$ & 93.7 млн gas \\
MNT6 full equation & два Miller loop + packed FE & 226.1 млн gas \\
MNT6 Article640 residue & общий multi-Miller accumulator & 172.0 млн gas \\
KZG opening BN254 & короткая on-chain opening-проверка & 133 тыс. gas \\
KZG + MNT4 relation & BN254 non-native sparse Miller estimate & 63.3 млн constraints \\
Merkle/ordinary-FRI & ожидаемая модельная стоимость & 78.3 млн gas \\
Merkle/ordinary-FRI & размер доказательства & 2.07 МБ \\
MNT-cycle-native relation & MNT4 + MNT6 fragments & 73,068 операций \\
lollipop-305 & полный оптимизированный fixed-shards цикл & 83.3 млн gas \\
\bottomrule
\end{longtable}

\subsection{Три формы стоимости}

Полученные результаты показывают не один, а три взаимосвязанных вида расходов.

\paragraph{On-chain arithmetic.}
Если считать MNT-сопряжение непосредственно в EVM, приходится платить gas за
трехсловную арифметику, расширения поля и сотни шагов цикла Миллера.

\paragraph{Off-chain constraints.}
Если перенести вычисление в proof над BN254, MNT-арифметика становится
non-native. Gas on-chain уменьшается, но растет число constraints, размер
witness и время работы prover-а.

\paragraph{Calldata и proof size.}
Если использовать нативные таблицы и Merkle/FRI, можно уменьшить зависимость от
BN254 non-native arithmetic, но стоимость возвращается через Merkle-пути,
раскрытия таблиц и calldata.

\subsection{Экономическая интерпретация}

На текущий момент на Etherscan:
\[
    \text{ETH price}\approx1982.85\text{ USD},
\]
\[
    \text{Ethereum standard gas price}\approx0.294\text{ gwei}.
\]
RPC Arbitrum One возвращал:
\[
    \text{Arbitrum One execution gas price}\approx0.020\text{ gwei}.
\]
Для Ethereum mainnet:
\[
    \text{cost}_{USD}=
    \Gas\cdot \text{gas price}_{gwei}\cdot10^{-9}\cdot1982.85.
\]

\begin{table}[htbp]
\centering
\caption{Стоимость эквивалентного EVM-исполнения по снимку 01.06.2026}
\small
\begin{tabularx}{\linewidth}{Y r r r}
\toprule
Режим & Gas & Ethereum mainnet & Arbitrum L2 execution \\
\midrule
MNT4 prepared sparse & 79,586,596 & \$46.40 & \$3.16 \\
MNT4 Article640 residue & 93,734,789 & \$54.64 & \$3.72 \\
MNT6 Article640 residue & 172,004,717 & \$100.27 & \$6.82 \\
lollipop full optimized cycle & 83,270,967 & \$48.54 & \$3.30 \\
Merkle/FRI lower bound & 51,359,352 & \$29.94 & \$2.04 \\
Merkle/FRI expected model & 78,340,624 & \$45.67 & \$3.11 \\
\bottomrule
\end{tabularx}
\end{table}

\section{Многошаговое исполнение тяжелого on-chain вычисления}

\subsection{Зачем требуется разбиение}

Многошаговый контракт не реализован в текущем коде. Ниже зафиксирована
архитектура следующего этапа, которая непосредственно опирается на
реализованные Miller hot path и fixed-cache режимы.

Даже оптимизированные verifier-ы остаются дороже лимита одной транзакции
Ethereum. Это не означает, что вычисление невозможно выполнить в EVM:
последовательный цикл Миллера можно разбить на несколько транзакций. Такой
режим не уменьшает суммарный gas и добавляет расходы на хранение состояния,
но позволяет выполнить длинное вычисление без превышения лимита отдельного
блока. Он полезен как промежуточный практический режим для L2, тестовых сетей
и специализированных приложений, где допустимы несколько последовательных
вызовов.

\subsection{Stateful checkpoint-контракт}

Для каждой проверки контракт создает сессию:
\[
    \mathsf{Session}=
    (\mathsf{statementHash},\mathsf{cacheId},
    i,f,\mathsf{active}).
\]
Здесь $\mathsf{statementHash}$ связывает сессию с публичными входами,
$\mathsf{cacheId}$ задает версию фиксированного кэша линий, $i$ является
номером следующего шага Миллера, а $f$ --- сохраненным аккумулятором.
Для $\widehat E$ вместо одного значения сохраняются два:
\[
    f_{\mathrm{num}},\qquad f_{\mathrm{den}}.
\]

Минимальный API имеет вид:
\begin{verbatim}
startSession(statement, cacheId) -> sessionId
advanceMiller(sessionId, maxSteps)
finalizeResidue(sessionId, c) -> bool
\end{verbatim}

\code{startSession} проверяет публичные точки, фиксирует хеш утверждения и
инициализирует аккумулятор единицей. \code{advanceMiller} обрабатывает
не более \code{maxSteps} последовательных строк кэша:
\[
    f\leftarrow f^2\ell_i(P)
\]
для doubling-шагов и
\[
    f\leftarrow f\ell_i(P)
\]
для addition-шагов. Номер шага берется из состояния контракта и после
успешного вызова увеличивается внутри контракта. Пользователь не может
пропустить шаг или заменить промежуточный аккумулятор.

\code{finalizeResidue} разрешен только при $i=N$, где $N$ --- точное число
шагов выбранного цикла. Он выполняет финальную проверку с
$c$-свидетельством, возвращает результат и удаляет состояние завершенной
сессии. Для защиты от накопления брошенных сессий можно использовать
депозит, срок жизни и отдельную функцию очистки после истечения срока.

\subsection{Граница применимости}

Адреса code-shards должны фиксироваться при развертывании или однозначно
определяться выбранным $\mathsf{cacheId}$. Состояние сессии нельзя принимать
из calldata без дополнительного доказательства: иначе пользователь мог бы
подменить аккумулятор. Stateful checkpoint-контракт решает именно проблему
лимита одной транзакции. Для снижения \emph{суммарного} gas требуется
дополнительное сжатие трассы, например polynomial relation check, KZG или
Merkle/FRI.

\section{Практическая полезность и дальнейшее развитие}

\subsection{Что можно использовать уже сейчас}

Работа не дает готовый дешевый однотранзакционный MNT-verifier для Ethereum
mainnet, однако ее результаты применимы на практике:
\begin{enumerate}
    \item библиотека Solidity/Yul-арифметики является воспроизводимым
    ориентиром стоимости MNT4-753, MNT6-753 и двухсловных полей;
    \item fixed-cache verifier-ы позволяют оценивать и прототипировать
    приложения с заранее известными $G_2$-точками;
    \item предложенная stateful multistep-архитектура показывает, как довести
    тяжелое вычисление до исполнения в EVM, если приложение допускает
    несколько транзакций;
    \item gas-таблицы и нижние оценки позволяют заранее исключать архитектуры,
    которые не помещаются в бюджет конкретной сети;
    \item cycle-native accounting дает основу для реализации следующего
    этапа: folding-слоя, где арифметика одной кривой выражается в поле другой
    нативно.
\end{enumerate}

Наиболее реалистичный прикладной сценарий текущего кода --- исследовательский
или специализированный L2-контур с фиксированными параметрами, code-shards и
многошаговым исполнением. Для публичного mainnet-verifier-а общего назначения
потребуется дальнейшее сжатие проверки.

\subsection{Направления дальнейших исследований}

\begin{enumerate}
    \item Реализовать stateful checkpoint-контракт и измерить накладные расходы
    хранения состояния для MNT4, MNT6 и lollipop.
    \item Реализовать MNT4/MNT6 cycle-native folding layer и проверить
    экспериментальную оценку constraints на полном рекурсивном контуре.
    \item Продолжить исследование polynomial relation check: сравнить KZG,
    Merkle/FRI и возможные гибридные схемы с учетом одновременно gas,
    constraints и размера calldata.
    \item Исследовать безопасные кривые меньшего размера по направлению
    lollipop и проверить, сохраняется ли выигрыш двухсловной арифметики при
    production-уровне стойкости.
    \item Рассмотреть специализированную поддержку сети: MNT-precompile,
    системный контракт или L2 с более подходящей моделью стоимости.
\end{enumerate}


\section{Итоговые выводы}

\begin{enumerate}
    \item Реализована и измерена оптимизированная 753-битная арифметика
    MNT4-753/MNT6-753 в Solidity/Yul. Для MNT4 экспериментально обоснован выбор
    Montgomery/CIOS, Karatsuba, специализированных умножений на нерезиденты,
    sparse lines и packed hot path.

    \item Построена оптимизационная лестница от наивного Tate baseline
    ($2.55$ млрд gas по строгой нижней экстраполяции) до prepared MNT4 пути
    ($79.6$ млн gas). Это показывает эффект оптимизаций и одновременно
    сохраняющуюся границу EVM.

    \item Реализована идея residue-проверки по ePrint 2024/640. Для MNT4
    финальная часть сокращается с $28.25$ млн до $6.13$ млн gas overhead.
    Для MNT6 общий multi-Miller residue verifier снижает equation baseline на
    $54.07$ млн gas.

    \item Исследована замена цикла Миллера полиномиальной проверкой. KZG
    уменьшает on-chain gas, но переносит стоимость в BN254 non-native
    constraints. Merkle/FRI избегает части этого переноса, но требует большого
    proof и calldata.

    \item Реализован предварительный MNT4/MNT6 cycle-native слой. Его
    accounting дает $24{,}126$ и $48{,}942$ multiplication operations для
    подготовленных fragments.

    \item Исследована конструкция lollipop-305. Двухсловная арифметика
    действительно дешевле MNT4 примерно в $2.8$--$3.3$ раза. Для третьей
    кривой реализован optimized Ehat path: product-accumulator для $P$ и $-P$,
    специализированный квадрат $\mathbb{F}_{q^6}$, scratch arena,
    pointer-swap, sparse multiplication и Frobenius-декомпозиция
    $c^p=c^q(c^{-1})^{q-p}$. Стоимость третьей кривой снижена до
    $56{,}163{,}308$ gas, а полный lollipop fixed-shards цикл --- до
    $83{,}270{,}967$ gas.

    \item Для вычислений, превышающих лимит отдельной транзакции, предложена
    stateful multistep-архитектура. Она сохраняет аккумулятор Миллера между
    вызовами и позволяет выполнить проверку по частям, не выдавая снижение
    суммарного gas за протокольное сжатие.
\end{enumerate}

Главный вывод работы:
\begin{quote}
Без MNT-precompile или нового proof-system слоя нельзя одновременно получить
малую on-chain стоимость и малую off-chain стоимость простым переносом
вычислений. Выполненная работа локализует причины этого ограничения,
подтверждает их кодом и измерениями и показывает наиболее перспективное
направление продолжения: MNT4/MNT6 cycle-native relation layer с последующим
folding.
\end{quote}

\subsection{Закрытие замечаний предыдущего отчета}

\begin{longtable}{p{0.28\linewidth}p{0.63\linewidth}}
\\
\toprule
Замечание & Что выполнено \\
\midrule
\endfirsthead
\toprule
Замечание & Что выполнено \\
\midrule
\endhead
Сравнить алгоритмы умножения и редукции &
Реализованы и измерены Montgomery/CIOS, Barrett, FIOS, Comba/SOS,
square-Comba, lazy reduction и branchless reduction. Для башен расширений
исследованы Karatsuba и специальные умножения. \\
\addlinespace
Обосновать оптимальность реализации &
Построена opcode-only нижняя граница 3-limb умножения, реализован ручной
all-stack Yul hot path, исследованы pointer/packed API, scratch arena,
fused sparse multiplication и потоковое чтение code-shards. После fused
оптимизации локальный резерв улучшения составляет доли процента полного
Miller path. \\
\addlinespace
Дать оценки сложности &
Оценки построены на нескольких уровнях: EVM opcode, $F_q$,
$F_{q^2}/F_{q^4}$ и $F_{q^3}/F_{q^6}$ операции, число Miller rounds,
стоимость FE, calldata, Merkle paths и constraints.  \\
\addlinespace
Исследовать более совершенные техники &
Реализованы shared multi-Miller accumulator, prepared sparse lines,
residue-проверка FE, KZG-opening, Merkle/FRI эксперимент, FRI cost model, MNT4/MNT6 cycle-native accounting и lollipop-305 по ePrint 2024/1627. \\
\bottomrule
\end{longtable}

\begin{thebibliography}{9}
\bibitem{eprint640}
Andrija Novakovic, Liam Eagen.
\emph{On Proving Pairings}.
Cryptology ePrint Archive, Paper 2024/640, 2024.
URL: \url{https://eprint.iacr.org/2024/640.pdf}.

\bibitem{eprint1627}
Craig Costello, Pratyush Korpal.
\emph{Cycles of supersingular elliptic curves for pairing-based proof systems}.
Cryptology ePrint Archive, Paper 2024/1627, 2024.
URL: \url{https://eprint.iacr.org/2024/1627.pdf}.

\bibitem{sonobe}
Privacy and Scaling Explorations.
\emph{Sonobe: a modular library for folding schemes}.
URL: \url{https://github.com/privacy-ethereum/sonobe}.

\bibitem{etherscanGas}
Etherscan.
\emph{Ethereum Gas Tracker}.
URL: \url{https://etherscan.io/gastracker}.
Снимок данных: 01.06.2026, 13:51 MSK.

\bibitem{arbitrumRpc}
Offchain Labs.
\emph{Arbitrum One public RPC endpoint}.
URL: \url{https://arb1.arbitrum.io/rpc}.
Снимок данных: 01.06.2026, 13:51 MSK.

\bibitem{arbitrumFees}
Offchain Labs.
\emph{Arbitrum Nitro: gas and fees}.
URL: \url{https://docs.arbitrum.io/how-arbitrum-works/gas-fees}.
\end{thebibliography}

\end{document}
